%%%PAGE 236 rot=0 legibility=good summary=摩擦力：滑动摩擦与静摩擦的特点、斜面上静摩擦方向判断、摩擦力突变（斜面f-t图、F=kt靠墙、小车上物块）
%%%BLOCK id=236-01 ch=02 sec=滑动摩擦力 kind=knowledge alt=none cont=none y=0.05-0.33
\textbf{滑动摩擦}　产生条件有弹力　$\mu N$

阻碍相对运动

方向与相对运动方向相反，与运动方向不一定相同

$f_{\text{滑}}$不一定是阻力

%FIG p236_a 0.145 0.158 0.275 0.225
\notefig[0.25]{figures/p236_a.png}

%FIG p236_b 0.21 0.208 0.615 0.325
\notefig[0.4]{figures/p236_b.png}
%%%END
%%%BLOCK id=236-02 ch=02 sec=静摩擦力 kind=knowledge alt=none cont=none y=0.34-0.54
\textbf{静摩擦}　$0<f\le f_{\max}$ % CHECK: 手写左端像 "<"（物理上应为 $\le$），照录

$\left\{\begin{array}{l}N\\ \mu\end{array}\right.$

阻碍　相对运动趋势

与运动方向无关

与相对运动趋势方向相反

%FIG p236_c 0.372 0.36 0.68 0.424
\notefig[0.3]{figures/p236_c.png}

\[
\left\{\begin{array}{ll}
F>mg\sin\alpha & \text{下}\\
F=mg\sin\alpha & \text{无}\\
F<mg\sin\alpha & \text{上}
\end{array}\right.
\]

$f_{\text{静}}$方向：
\[
a\left\{\begin{array}{ll}
>g\sin\theta & \text{沿下}\\
=g\sin\theta & \text{无}\\
<g\sin\theta & \text{沿上}
\end{array}\right.
\]
%%%END
%%%BLOCK id=236-03 ch=02 sec=摩擦力突变 kind=method alt=none cont=none y=0.55-0.68
%FIG p236_d 0.10 0.558 0.2445 0.622
\notefig[0.25]{figures/p236_d.png}

$f_{\text{静}}=mg\sin\omega t$　　$f_{\text{滑}}=\mu mg\cos\omega t$

%FIG p236_e 0.225 0.612 0.44 0.675
\notefig[0.25]{figures/p236_e.png}
%%%END
%%%BLOCK id=236-04 ch=02 sec=摩擦力突变 kind=method alt=03 cont=none y=0.69-0.88
% 图中标注：F=kt，μ，f，G，f=μkt，向下减速，静止，向下加速，t
%FIG p236_f 0.045 0.695 0.34 0.875
\notefig[0.3]{figures/p236_f.png}
%%%END
%%%BLOCK id=236-05 ch=02 sec=摩擦力突变 kind=problem alt=03 cont=none y=0.88-1.0
%FIG p236_g 0.11 0.885 0.265 0.995
\notefig[0.25]{figures/p236_g.png}

$F_{\text{拉}}=5\unit{N}$　　$m=5\unit{kg}$　均静止　右加速 $a$ 从 $0\to 2\unit{m/s^2}$ % CHECK: "5" 与 F拉=5N 的手写数字同形，由 $a_1=F/m=1\unit{m/s^2}$ 判断 m=5kg；"均"为插入字

$F+f_{\max}=ma_2$　　$a_2\ge 2\unit{m/s^2}$　　$\therefore F_{\text{拉}}$不变，$f$先向左$\downarrow$ 再向右$\uparrow$

$f=F=5\unit{N}\le f_{\max}$　　$a_1=1\unit{m/s^2}$
%%%END
%%%PAGEEND 236
